% part: intuitionistic-logic
% chapter: semantics
% section: introduction

\documentclass[../../../include/open-logic-section]{subfiles}

\begin{document}

\olfileid{int}{sem}{int}

\olsection{प्रस्तावना}

चिन्हार्थमीमांसेशिवाय कोणत्याही तर्कशास्त्राचे
समाधानकारक वर्णन होत नाही; अंतःप्रज्ञावादी
तर्कशास्त्रही त्याला अपवाद नाही. अभिजात
तर्कशास्त्रात !!{valuation}s वर आधारलेली
चिन्हार्थमीमांसा प्रमाणभूत मानली जाते. पण
अंतःप्रज्ञावादी तर्कशास्त्रासाठी अनेक स्पर्धक
चिन्हार्थमीमांसा आहेत. संयोजकांच्या अर्थाचे
अंतःप्रज्ञावादी दृष्टिकोनातून मान्य होईल असे
पूर्ण स्पष्टीकरण देण्याच्या बाबतीत त्यांपैकी
एकही पूर्णपणे समाधानकारक नाही.

मोडल तर्कशास्त्रातील चिन्हार्थमीमांसेसारखी,
!!{relational model}s वर आधारलेली चिन्हार्थमीमांसा
बहुधा सर्वाधिक प्रचलित आहे. यात !!{propositional variable}s
जगांशी संबद्ध केली जातात आणि जगांना प्राप्यता
संबंध जोडतो. हा संबंध नेहमीच अंशतः क्रम असतो;
म्हणजे तो परावर्ती, प्रतिसममित आणि संक्रमक असतो.

सहज कल्पना करण्यासाठी या जगांना ज्ञानावस्था किंवा
“पुराव्याची परिस्थिती” समजू शकता. $w'$ ही अवस्था
$w$ पासून प्राप्य असेल तेव्हा आणि केवळ तेव्हाच,
आपल्याला सध्या माहीत असलेल्या गोष्टींनुसार $w'$
ही ज्ञानाची संभाव्य (भावी) अवस्था असते; म्हणजे
$w$ मध्ये जे माहीत आहे त्याच्याशी ती सुसंगत असते.
एकदा एखादे विधान माहीत झाले की ते नंतर अज्ञात
होऊ शकत नाही: $!A$ हे~$w$ मध्ये माहीत असेल
आणि~$Rww'$ असेल, तर $!A$ हे~$w'$ मध्येही माहीत
असते. त्यामुळे प्राप्यता संबंधाच्या दृष्टीने
“ज्ञान” एकस्वनिक आहे.

ज्ञानविषयक तर्कशास्त्राप्रमाणे “$!A$ माहीत आहे”
याचा अर्थ “सर्व ज्ञानविषयक पर्यायांत सत्य आहे”
असा घेतला, तर $!A \land !B$ हे~$w$ मध्ये माहीत
असण्याची अट अशी आहे: सर्व ज्ञानविषयक पर्यायांत
$!A$ आणि~$!B$ दोन्ही माहीत असावेत. पण ज्ञान
एकस्वनिक आहे आणि $R$ परावर्ती आहे. म्हणून
$!A \land !B$ हे~$w$ मध्ये माहीत असते तेव्हा
आणि केवळ तेव्हाच $!A$ आणि $!B$ दोन्ही~$w$
मध्ये माहीत असतात. याच कारणाने $!A \lor !B$
हे~$w$ मध्ये माहीत असते तेव्हा आणि केवळ तेव्हाच
त्यांपैकी किमान एक तेथे माहीत असते.
म्हणून $\land$ आणि $\lor$ यांच्या सत्यतेच्या अटी
अभिजात तर्कशास्त्रातील अटींशी जुळतात.

सशर्त विधानाच्या सत्यतेच्या अटी मात्र अभिजात
तर्कशास्त्राहून वेगळ्या आहेत. $!A \lif !B$
हे~$w$ मध्ये माहीत असते तेव्हा आणि केवळ तेव्हाच
$Rww'$ असलेली अशी कोणतीही~$w'$ अवस्था नसते
जिथे $!A$ माहीत आहे पण $!B$ माहीत नाही.
ही अट, $w$ मध्ये $!A$ माहीत नाही किंवा
$!B$ माहीत आहे, या अटीसारखी नाही. कारण
$w$ मध्ये $!A$ किंवा $!B$ यांपैकी काहीही माहीत
नसेल, तरी $Rww'$ असलेली अशी भावी ज्ञानावस्था
$w'$ असू शकते; त्या~$w'$ मध्ये $!A$ माहीत होते
पण $!B$ माहीत होत नाही.

ज्या कोणत्याही संभाव्य भावी ज्ञानावस्थेत~$!A$
माहीत होते अशी अवस्था नसेल, तरच आपल्याला
$\lnot !A$ माहीत असते. यामागील कल्पना अशी:
$!A$ जाणून घेणे शक्य असेल, तर एखाद्या
संभाव्य भावी ज्ञानावस्थेत~$!A$ माहीत होते.
पण~$\lfalse$ माहीत असणे शक्य नाही; म्हणून
त्या भावी ज्ञानावस्थेत आपल्याला $!A$ माहीत
असेल, पण~$\lfalse$ माहीत नसेल.

या अर्थलक्षणात वर्जित मध्याचे तत्त्व वैध ठरत नाही.
कारण अशी काही $!A$ सूत्रे आहेत की जी अद्याप
माहीत नाहीत, पण पुढे माहीत होऊ शकतात.
अशा !!a{formula}~$!A$ साठी $!A$ आणि~$\lnot !A$
दोन्ही माहीत नसतात; म्हणून $!A \lor \lnot !A$
माहीत नसते. पण उदाहरणार्थ, $\lnot(!A
\land \lnot !A)$ हे आपल्याला माहीत असते.
कारण $!A$ आणि~$\lnot !A$ दोन्ही माहीत असलेली
कोणतीही भावी अवस्था शक्य नसते; आणि हे
$!A$ किंवा~$\lnot !A$ माहीत आहे की नाही
यापासून स्वतंत्रपणे आपल्याला माहीत असते.

अंतःप्रज्ञावादी तर्कशास्त्रासाठी !!^{relational model}s
हीच एकमेव उपलब्ध चिन्हार्थमीमांसा नाही.
सांस्थितिक चिन्हार्थमीमांसा हा आणखी एक पर्याय आहे:
त्यात विधानांचे अर्थलक्षण सांस्थितिक अवकाशातील
विवृत संच म्हणून केले जाते आणि संयोजकांचे
अर्थलक्षण त्या संचांवरील क्रिया म्हणून केले जाते
(उदा., $\land$ ला छेद अनुरूप असतो).

\end{document}
